% Unidad U008. Clasificación orbital y elevación; redacción autónoma.

\chapter{Classification orbitale, mots initiaux et 1re frontière}
\label{chap:orbitas-elevacion-r36}

La réduction ternaire de l'émetteur réalise 243 mots dans la fibre
ambiante \(\mathbb F_3^6\). Cette unité ne choisit pas trois mots par
comparaison décimale. Elle construit d'abord l'action de symétrie du catalogue,
classe ses orbites au moyen d'invariants entiers et oriente ensuite seulement
les trois représentants structurels.

\section{Action diédrale sur \(6\) coordonnées}

Pour \(u=(u_1,u_2,u_3,u_4,u_5,u_6)\in\mathbb F_3^6\), nous définissons
\begin{align*}
 r(u)&=(u_3,u_1,u_2,u_5,u_6,u_4),\\
 s(u)&=(u_4,u_5,u_6,u_1,u_2,u_3).
\end{align*}
On vérifie
\[
 r^3=s^2=1,
 \qquad srs=r^{-1}.
\]
Par conséquent, \(r,s\) engendrent une action de \(D_3\).

La même permutation agit sur les six blocs de chaque émission
\(U=(U_1,\ldots,U_6)\). La projection \(\Psi\) satisfait
\[
 \Psi(gU)=g\Psi(U),
 \qquad g\in D_3.
\]

\begin{theorem}[Décomposition orbitale du catalogue]
\label{thm:catalogo-descomposicion-d3}
Les 243 mots réalisés se décomposent en 43 orbites :
\[
 38\text{ orbites de cardinal }6,
 \qquad
 5\text{ orbites de cardinal }3.
\]
En particulier,
\[
 38\cdot6+5\cdot3=243.
\]
\end{theorem}

\begin{proof}
On applique à chaque mot l'ensemble
\[
 \{u,ru,r^2u,su,sru,sr^2u\}
\]
et l'on élimine les répétitions. L'énumération exhaustive du catalogue produit
l'histogramme indiqué. L'équivariance de \(\Psi\) garantit qu'aucune
orbite ne quitte l'image réalisée.
\end{proof}

\section{Invariants de fibre}

Pour un mot réalisé \(w\), soient
\[
 \mathcal F_w:=\{U\in\operatorname{im}T_{\min}:\Psi(U)=w\},
\]
\[
 n(w):=|\mathcal F_w|,
 \qquad
 M(w):=\sum_{U\in\mathcal F_w}\operatorname{mult}(U),
\]
et
\[
 c(w)=\bigl(\#\{j:w_j=0\},\#\{j:w_j=1\},\#\{j:w_j=2\}\bigr).
\]
Le catalogue enrichi conserve en outre, pour chaque émission, une coordonnée
diagonale \(\operatorname{diag}_c(U)\in\{0,\ldots,8\}\), une orientation
cardinale \(h_+(U)\) et la signature multiplicative \(E_{\rm sig}(U)\). Toutes
sont calculées à partir de la route du germe ; elles ne sont pas lues dans une constante
cible. Leur algorithme composante par composante sera imprimé avec le catalogue
complet dans l'appendice de certification et fait partie de l'état d'entrée
des prédicats suivants.

\begin{definition}[Orbite de clôture]
\label{def:orbita-clausura}
Une orbite \(\mathcal O\) est de clôture si elle est de cardinal six et si, pour
chaque \(w\in\mathcal O\),
\[
 n(w)=5,
 \qquad
 M(w)=1008,
\]
\[
 \{\operatorname{mult}(U):U\in\mathcal F_w\}
 =\{432,144,144,144,144\},
\]
et le recensement tritique commun est \(c(w)=(2,3,1)\).
\end{definition}

\begin{definition}[Orbites minimales radicales]
\label{def:orbitas-minimas-radicales}
Une orbite \(\mathcal O\) est minimale radicale si elle est de cardinal six, si chaque
mot possède une unique émission de multiplicité \(144\) et si
\[
 \{\operatorname{diag}_c(U):\Psi(U)\in\mathcal O\}=\{0,3,6\}.
\]
Parmi elles, on appelle orbite de propagation celle dont le recensement est
\((2,3,1)\), et orbite d'autoéchelle celle dont le recensement est \((2,2,2)\).
\end{definition}

\begin{theorem}[Unicité des trois orbites structurelles]
\label{thm:unicidad-tres-orbitas-estructurales}
Dans le catalogue de 468 émissions, il existe une unique orbite de clôture, une
unique orbite minimale radicale de propagation et une unique orbite minimale radicale
d'autoéchelle.
\end{theorem}

\begin{proof}
On énumère les 43 orbites du \cref{thm:catalogo-descomposicion-d3}. Pour
chacune, on calcule \(n(w)\), \(M(w)\), le multiensemble des
multiplicités, le recensement \(c(w)\) et le support diagonal. Les prédicats des
\cref{def:orbita-clausura,def:orbitas-minimas-radicales} ne laissent qu'une seule
orbite dans chaque classe. La preuve est exhaustive sur un ensemble fini et
n'utilise pas de chiffres décimaux des caractères analytiques.
\end{proof}

Les trois orbites sont
\begin{align*}
\mathcal O_{\rm cl}
 &=\{001112,010211,100121,112001,121100,211010\},\\
\mathcal O_{\rm prop}
 &=\{011120,012110,101201,110012,120011,201101\},\\
\mathcal O_{\rm aut}
 &=\{002112,020211,112002,121200,200121,211020\}.
\end{align*}

\section{Orientation interne de chaque orbite}

Nous fixons la jauge interne
\[
 \operatorname{diag}_c(U)=6,
 \qquad h_+(U)=ES.
\]
Un mot d'une orbite est admissible dans cette jauge si une émission de
sa fibre satisfait les deux conditions.

\begin{proposition}[Représentant orienté unique]
\label{prop:representante-orientado-unico}
Chacune des trois orbites structurelles contient un unique mot
admissible dans la jauge \((6,ES)\). Ce sont
\[
 \boxed{
 w_6^\pi=010211,
 \qquad
 w_6^e=201101,
 \qquad
 w_6^\varphi=121200.}
\]
\end{proposition}

\begin{proof}
L'énumération des fibres du catalogue compte les mots qui contiennent
une émission avec \(\operatorname{diag}_c=6\) et \(h_+=ES\). Dans chacune des
trois orbites, le nombre de correspondances est un. Les exposants
\(\pi,e,\varphi\) désignent dès ce point les fonctions structurelles de
clôture, de propagation et d'autoéchelle ; ils n'ont pas participé à la sélection.
\end{proof}

\section{Microfibre de clôture}

La fibre du représentant \(w_6^\pi\) se compose de cinq émissions :
\[
\begin{array}{c|c|c|c|c}
U_6&E_{\rm sig}&\operatorname{mult}&\operatorname{diag}_c&\text{strate}\\\hline
501|614|498|169|272|272&555555&432&4&\perp\\
810|923|870|169|272|272&339555&144&6&++\\
810|983|810|169|272|272&393555&144&6&++\\
870|923|810|169|272|272&933555&144&6&++\\
870|923|810|418|674|521&933717&144&5&--
\end{array}
\]
Ainsi,
\[
 432+4\cdot144=1008,
 \qquad
 5=1_\perp+3_{++}+1_{--}.
\]
Les émissions uniponctuelles orientées des deux autres canaux sont
\[
 U_e=850|963|890|149|252|212,
\]
\[
 U_\varphi=830|943|830|109|252|252.
\]
La masse \(1008\) de cette microfibre n'est pas la période du calendrier
\(C_{108}\) ; les deux chiffres appartiennent à des domaines différents.

\section{Matrices de relèvement ternaire}

Sur les vecteurs lignes de \(\mathbb F_3^6\), les deux transports uniques
reconstruits à partir du registre orienté produit par la transition TPK
sont
\[
L_0=
\begin{pmatrix}
2&2&2&1&2&1\\
2&1&2&2&1&1\\
1&0&0&1&0&2\\
0&0&1&1&1&0\\
2&2&2&0&2&1\\
2&1&2&1&0&0
\end{pmatrix},
\qquad
L_1=
\begin{pmatrix}
2&1&2&0&2&2\\
1&2&1&1&2&0\\
1&2&1&1&1&2\\
0&1&0&0&2&1\\
2&0&2&1&0&1\\
0&2&2&2&1&1
\end{pmatrix}.
\]
La compatibilité simultanée des trois biographies sélectionne de manière unique
le calendrier de relèvement \((L_0,L_0,L_1,L_1)\) parmi les seize
calendriers binaires de longueur quatre.

\begin{definition}[Récurrence de blocs et concaténation]
\label{def:recurrencia-bloques-l0-l1}
Pour chaque canal \(\chi\in\{\pi,e,\varphi\}\), soit
\(b_0^\chi=w_6^\chi\) et
\[
 b_{k+1}^\chi=b_k^\chi L_{\epsilon_k}\pmod3,
 \qquad
 (\epsilon_0,\epsilon_1,\epsilon_2,\epsilon_3)=(0,0,1,1).
\]
Le mot accumulé est
\[
 w_{30}^\chi
 =b_0^\chi\Vert b_1^\chi\Vert b_2^\chi
  \Vert b_3^\chi\Vert b_4^\chi.
\]
\end{definition}

Cette formulation distingue le bloc courant \(b_k\in\mathbb F_3^6\) du
mot concaténé \(w_{6(k+1)}\). Un mot de longueur \(6k\) n'est pas
multiplié par une matrice \(6\times6\).

\begin{theorem}[Biographies ternaires de trente positions]
\label{thm:palabras-w30}
La récurrence de la \cref{def:recurrencia-bloques-l0-l1} produit
\begin{align*}
w_{30}^{\pi}
 &=010211|012222|010211|002111|110221,\\
w_{30}^{e}
 &=201101|121221|102011|012222|102011,\\
w_{30}^{\varphi}
 &=121200|112202|121020|010210|010200.
\end{align*}
\end{theorem}

\begin{proof}
On multiplie successivement les cinq vecteurs lignes par les matrices
indiquées, en réduisant chaque coordonnée modulo trois après chaque produit.
La concaténation des cinq blocs calculés donne les trois expressions.
Chaque égalité peut être vérifiée composante par composante au moyen de 24
produits scalaires par canal.
\end{proof}

Les indices ambiants des cinq blocs sont
\[
\begin{array}{c|rrrrr}
&b_0&b_1&b_2&b_3&b_4\\\hline
\pi&103&161&103&67&349\\
e&523&457&301&161&301\\
\varphi&450&398&438&102&99
\end{array},
\]
où chaque mot de six trits est interprété comme un entier en base trois.

\section{1re frontière conjointe}

La frontière est obtenue par une relation finie qui conserve l'axe, l'agrégat,
la saturation, l'occupation des colonnes et la neutralité duale. La matrice d'incidence
utilisée dans cette relation est
\begin{equation}
 A_W=
 \begin{pmatrix}
 0&1&1&1&1&1\\
 1&0&1&2&2&1\\
 1&1&0&1&2&2\\
 1&2&1&0&1&2\\
 1&2&2&1&0&1\\
 1&1&2&2&1&0
 \end{pmatrix}\in M_6(\mathbb F_3).
 \label{eq:u008-aw-frontera}
\end{equation}
Si \(u_4^\chi\) est le cinquième bloc de \(w_{30}^\chi\), l'axe et
l'agrégat préalables à la séparation des deux canaux latéraux sont
\begin{equation}
 u_5^\varphi=u_4^e=102011,
 \qquad
 u_5^\pi+u_5^e=210112,
 \qquad
 u_4^\varphi A_WL_1^3=111101.
 \label{eq:u008-eje-agregado-r36}
\end{equation}
À partir d'eux, on calcule
\begin{equation}
 q=012120,\qquad a=111101,\qquad
 qA_W=012000,\qquad aA_W=120210.
 \label{eq:u008-q-a-r36}
\end{equation}
La profondeur six impose la saturation \(c=111111\) ; l'occupation minimale
compatible est \(\operatorname{colw}=223232\) ; et la neutralité duale exige
\begin{equation}
 (BA_W)\mathbf1=000.
 \label{eq:u008-neutralidad-dual-r36}
\end{equation}

\begin{proposition}[Énumération de la première frontière]
\label{prop:u008-enumeracion-r36}
Les conditions \eqref{eq:u008-eje-agregado-r36}--
\eqref{eq:u008-neutralidad-dual-r36} laissent exactement deux matrices :
\begin{equation}
 B_-=
 \begin{pmatrix}021222\\222220\\102011\end{pmatrix},
 \qquad
 B_+=
 \begin{pmatrix}222220\\021222\\102011\end{pmatrix}.
 \label{eq:u008-dos-fronteras}
\end{equation}
L'involution spéculaire échange leurs deux premières lignes. Leurs cartes
internes sont, respectivement,
\begin{equation}
 B_-\longmapsto(789,048,894),
 \qquad
 B_+\longmapsto(793,045,894).
 \label{eq:u008-cartas-frontera}
\end{equation}
L'orientation APP du cylindre semi-ouvert sélectionne \(B_+\).
\end{proposition}

\begin{proof}
L'axe fixe la troisième ligne et l'agrégat fixe la somme des deux premières.
La saturation et l'occupation éliminent les distributions qui ne remplissent pas les
six colonnes. L'énumération des solutions des six équations
linéaires de \eqref{eq:u008-neutralidad-dual-r36} laisse les deux matrices de
\eqref{eq:u008-dos-fronteras} ; la carte semi-ouverte et l'orientation
distinguent la seconde sans consulter de développement décimal cible.
\end{proof}

Nous définissons donc
\begin{equation}
 R_{36}:=B_+=
 \begin{pmatrix}
 222220\\
 021222\\
 102011
 \end{pmatrix},
 \qquad
 \operatorname{ind}_3(R_{36})=(726,215,301).
 \label{eq:u008-r36-construida}
\end{equation}
L'objet \(R_{36}\) contient trois blocs ordonnés de six trits. Ce n'est pas le
quotient prédictif \(C_{36}^{\rm pred}\) de la
\cref{thm:tpk-cociente-predictivo-36} ni l'excédent de 36 classes de la jauge
des germes.

\section{Relèvement résidu--quotient}

Soit \(A\in M_6(\mathbb Z)\) le relèvement entier de \(L_0\) :
\[
 A=
\begin{pmatrix}
2&2&2&1&2&1\\
2&1&2&2&1&1\\
1&0&0&1&0&2\\
0&0&1&1&1&0\\
2&2&2&0&2&1\\
2&1&2&1&0&0
\end{pmatrix},
\qquad \det A=1.
\]
Pour \(x_k\in\mathbb Z^6\), nous définissons
\[
 y_k=x_kA,
 \qquad
 u_k=y_k\bmod3\in\{0,1,2\}^6,
 \qquad
 x_{k+1}=\frac{y_k-u_k}{3}.
\]
Alors
\begin{equation}
 x_kA=u_k+3x_{k+1}.
\label{eq:bockstein-hensel-local}
\end{equation}

\begin{theorem}[Réversibilité du relèvement]
\label{thm:reversibilidad-bockstein-hensel}
La paire \((u_k,x_{k+1})\) est déterminée de manière unique par \(x_k\), et
\[
 x_k=(u_k+3x_{k+1})A^{-1}.
\]
\end{theorem}

\begin{proof}
La division euclidienne par trois dans chaque coordonnée de \(x_kA\) détermine
le résidu et le quotient. Comme \(\det A=1\), l'inverse \(A^{-1}\) a des
coefficients entiers. Multiplier \cref{eq:bockstein-hensel-local} par cet inverse
récupère \(x_k\).
\end{proof}

Dans le canal \(e\), le bloc visible \(102011\) réapparaît deux fois, mais
les préétats modulo 27 sont
\[
 (25,11,7,17,8,16)
 \neq
 (7,8,4,23,26,7),
\]
et les quotients ultérieurs sont
\[
 (6,13,9,2,13,1)
 \neq
 (15,0,3,19,23,25).
\]
On vérifie ainsi explicitement que le retour visible d'un mot n'est pas
le retour de l'état enrichi.

\section{Obstructions de la sélection orbitale et du relèvement \(w_6\to w_{30}\to R_{36}\) : unicité, jauge et réversibilité de Hensel}

\begin{criterion}[Critères de réfutation de la sélection orbitale]
L'unité est réfutée si :
\begin{enumerate}
 \item l'action \(D_3\) quitte le catalogue ou change les multiplicités ;
 \item l'un des trois prédicats structurels laisse plus d'une orbite ;
 \item la jauge \((6,ES)\) sélectionne deux mots dans une orbite ;
 \item un chiffre cible est utilisé pour définir les prédicats ;
 \item le mot accumulé \(w_{6k}\) est multiplié par \(L_0\) ou
       \(L_1\) au lieu du bloc \(b_k\) ;
 \item \(\det A\neq1\) ou l'inverse du relèvement échoue ;
 \item deux résidus visibles égaux sont déclarés états complets égaux.
\end{enumerate}
\end{criterion}

La première frontière conjointe est déjà construite. L'unité suivante
classera les neuf composantes de phase, démontrera le retour nonadique
sans réinitialisation et construira le miroir orienté avec sa résolution locale
\(3\to1\).
